\documentclass[aspectratio=169]{beamer}
\input{header}
\coursecode{JKSSB}

\title{Excel Structure}
\subtitle{Lesson 24 --- Workbook, Worksheet and Cell Anatomy}
\author{JKSSB Computer Awareness}
\date{\today}

\begin{document}
\frame{\titlepage}

\begin{frame}{Excel and the Workbook}
  \begin{itemize}
    \item \key{Microsoft Excel} is a spreadsheet program that stores,
      calculates, and analyses data in a grid.
    \item A \key{workbook} is the Excel \key{file} that you save and open.
    \item One workbook can hold \key{many worksheets}, also called sheets.
    \item Worksheet \key{tabs} at the bottom of the window let you move
      between them.
  \end{itemize}
  \begin{block}{The one idea}
    A workbook is the file; a worksheet is a single page inside that file.
  \end{block}
\end{frame}

\begin{frame}{Workbook vs Worksheet}
  \begin{center}
    \begin{tabular}{lll}
      \toprule
      \textbf{Point} & \textbf{Workbook} & \textbf{Worksheet} \\
      \midrule
      What it is & The Excel file & One sheet inside a workbook \\
      How many & One file & Many per workbook \\
      Saved as & A single file name & Part of that file \\
      Tab & --- & Tab at the bottom \\
      \bottomrule
    \end{tabular}
  \end{center}
  \vspace{0.25cm}
  \begin{alertblock}{Trap alert}
    A workbook is \key{not} the same as a worksheet. The worksheet is a part
    of the workbook, not the whole file.
  \end{alertblock}
\end{frame}

\begin{frame}{The Worksheet Grid: Rows and Columns}
  \begin{itemize}
    \item Every worksheet is a grid of \key{rows} and \key{columns}.
    \item \key{Rows} run horizontally and are named by \key{numbers} on the
      left (1, 2, 3, \dots).
    \item \key{Columns} run vertically and are named by \key{letters} across
      the top (A, B, C, \dots).
    \item Column labels go A--Z, then AA--AZ, then BA, and so on up to the
      last column.
  \end{itemize}
  \muted{Read the row number on the left edge and the column letter on the
  top edge of the sheet.}
\end{frame}

\begin{frame}{Grid Dimensions (Direct Recall)}
  \begin{center}
    \begin{tabular}{ll}
      \toprule
      \textbf{Property} & \textbf{Value in Excel 2007 and later} \\
      \midrule
      Total rows & 1,048,576 (numbered 1 to 1,048,576) \\
      Total columns & 16,384 (labelled A to XFD) \\
      First column & A \\
      Last column & XFD \\
      \bottomrule
    \end{tabular}
  \end{center}
  \vspace{0.3cm}
  \begin{alertblock}{Remember the two numbers}
    Rows: \key{1,048,576}. Columns: \key{16,384}. The last column is
    \key{XFD}.
  \end{alertblock}
\end{frame}

\begin{frame}{The Cell}
  \begin{itemize}
    \item A \key{cell} is the \key{intersection} of one row and one column.
    \item It is the basic box in which you type data.
    \item Every cell has a unique address, formed from its column letter and
      its row number.
    \item A worksheet is simply a very large collection of these boxes.
  \end{itemize}
  \begin{block}{Definition}
    A cell is the small rectangular box formed where a row and a column
    meet.
  \end{block}
\end{frame}

\begin{frame}{Cell Address (Cell Reference)}
  \begin{itemize}
    \item A \key{cell address} names one specific cell.
    \item It is written as the \key{column letter first}, then the
      \key{row number}.
    \item Example: \key{B5} means column B, row 5.
    \item Example: \key{C7} means column C, row 7.
  \end{itemize}
  \begin{exampleblock}{Worked example}
    The cell at column D and row 12 is written \key{D12} --- letter first,
    then number. Writing 12D or D:12 is wrong.
  \end{exampleblock}
\end{frame}

\begin{frame}{The Active Cell}
  \begin{itemize}
    \item The \key{active cell} is the cell that is currently
      \key{selected}.
    \item It is surrounded by a thick border.
    \item Whatever you type is entered into the active cell.
    \item Its address appears in the \key{Name Box}, and its contents appear
      in the \key{Formula Bar}.
  \end{itemize}
  \begin{alertblock}{Trap alert}
    Only \key{one} cell is active at a time, even when a whole range of cells
    is selected.
  \end{alertblock}
\end{frame}

\begin{frame}{The Formula Bar}
  \begin{itemize}
    \item The \key{Formula Bar} is the long white bar \key{below the Ribbon}
      and above the grid.
    \item It \key{shows the contents} of the active cell.
    \item For a formula, it shows the \key{formula itself}, not the result
      shown in the cell.
    \item You can also \key{edit} the active cell's contents directly in the
      Formula Bar.
  \end{itemize}
  \begin{block}{Keep it straight}
    The cell shows the \muted{result}; the Formula Bar shows the
    \key{formula or data} behind it.
  \end{block}
\end{frame}

\begin{frame}{The Name Box}
  \begin{itemize}
    \item The \key{Name Box} is the small box at the \key{left end of the
      Formula Bar}.
    \item It normally shows the \key{address of the active cell}.
    \item If a range is selected, it shows the \key{name or address of that
      range}.
    \item You can type an address into the Name Box to jump straight to that
      cell.
  \end{itemize}
\end{frame}

\begin{frame}{Range}
  \begin{itemize}
    \item A \key{range} is a group of two or more cells treated as one block.
    \item It is written with the two \key{corner addresses} separated by a
      colon (:).
    \item Example: \key{A1:C10} means columns A to C, rows 1 to 10.
    \item A range is usually \key{rectangular} and its cells are contiguous.
  \end{itemize}
  \begin{exampleblock}{Worked example}
    The range B2:D5 covers 3 columns (B, C, D) and 4 rows (2--5), so it
    contains $3 \times 4 = 12$ cells.
  \end{exampleblock}
\end{frame}

\begin{frame}{Cell vs Range (Near-Neighbour)}
  \begin{center}
    \begin{tabular}{lll}
      \toprule
      \textbf{Point} & \textbf{Cell} & \textbf{Range} \\
      \midrule
      What it is & One box & A block of cells \\
      Written as & B5 & B2:D5 \\
      Number of cells & Exactly one & Two or more \\
      \bottomrule
    \end{tabular}
  \end{center}
  \vspace{0.3cm}
  \begin{alertblock}{Trap alert}
    \key{B5} is a single cell, not a range. A range needs \key{two} corner
    addresses joined by a colon, such as \key{B5:D9}.
  \end{alertblock}
\end{frame}

\begin{frame}{Basic Data Types}
  \begin{itemize}
    \item \key{Text} (also called a label): words, names, headings.
    \item \key{Number}: values used in calculations.
    \item \key{Date and time}: stored as numbers but shown as dates.
    \item \key{Formula}: an entry beginning with \key{=} that calculates a
      result.
    \item Excel also accepts \key{TRUE/FALSE} (logical values) and shows
      errors such as \key{\#VALUE!}.
  \end{itemize}
\end{frame}

\begin{frame}{Text vs Numbers: Default Alignment}
  \begin{center}
    \begin{tabular}{lll}
      \toprule
      \textbf{Data type} & \textbf{Default alignment} & \textbf{Example} \\
      \midrule
      Text & Left & Roll No \\
      Number & Right & 4500 \\
      Date & Right & 15-08-2026 \\
      \bottomrule
    \end{tabular}
  \end{center}
  \vspace{0.3cm}
  \begin{block}{Quick check}
    If a value sits on the \key{right} of the cell, Excel has treated it as a
    \key{number}; if it sits on the \key{left}, it is \key{text}.
  \end{block}
\end{frame}

\begin{frame}{Moving Around a Worksheet}
  \begin{itemize}
    \item Click any cell to make it the \key{active cell}.
    \item Use the \key{arrow keys} or \key{Tab} / \key{Enter} to move.
    \item The worksheet \key{tabs} switch between sheets in a workbook.
    \item The \key{Name Box} jumps to a typed address; the Formula Bar reads
      and edits the active cell.
  \end{itemize}
  \muted{These four elements --- grid, Name Box, Formula Bar, tabs --- are the
  basic navigation anatomy of Excel.}
\end{frame}

\begin{frame}{Trap Alert: Facts to Keep Straight}
  \begin{itemize}
    \item A \key{workbook} is the file; a \key{worksheet} is one sheet inside
      it.
    \item A \key{cell} is one box; a \key{range} is two or more cells.
    \item Cell address = \key{column letter, then row number} (B5).
    \item Total rows: \key{1,048,576}; total columns: \key{16,384}.
    \item Last column: \key{XFD}.
    \item \key{Name Box} shows the address; \key{Formula Bar} shows the
      contents.
  \end{itemize}
\end{frame}

\begin{frame}{Lesson 24: Recall}
  \begin{itemize}
    \item A workbook is the Excel file; it can contain many worksheets.
    \item A worksheet is a grid of numbered rows and lettered columns.
    \item Excel 2007+ has 1,048,576 rows and 16,384 columns (A to XFD).
    \item A cell is the intersection of a row and a column; the active cell
      is the one selected.
    \item A cell address is column letter plus row number (B5); a range is a
      block such as A1:C10.
    \item The Name Box shows the address; the Formula Bar shows the contents.
    \item Basic data types: text, number, date/time, and formula.
  \end{itemize}
\end{frame}
\end{document}
